% !TeX root = ../Thesis.tex

%************************************************
\chapter{Introduction}\label{ch:introduction}
%************************************************
%***************************************************************************
\section{Motivation}
\label{sec:background_motivation}
%%\glsresetall % Resets all acronyms to not used
Data management and processing are essential aspects of many applications and digital services. The increasing volume of data and demand for quick and reliable data access have led to the need for more efficient and scalable databases \cite{rahimi2010distributed, ozsu1999principles, darmont2022advances}.  The advent of big data, cloud computing, and Internet of Things (IoT) has further intensified this need, as data is generated from a multitude of sources at an unprecedented rate. Traditional centralized databases, operating on a single machine, possess implicit limitations in handling the ever-growing data volume and the massive influx of user requests. This is due to their restricted processing and storage capacities. Although it is possible to scale up a single database server to a certain degree, an alternative approach of scaling out and leveraging multiple database nodes has become more prominent in recent trends \cite{ozsu1999principles}. With the capacity to distribute data over numerous database nodes, distributed databases have emerged as a preferred solution to limitations of centralized databases, resulting in enhanced speed, fault tolerance, and scalability \cite{silberschatz2011database, tomar2014overview, abdelhafiz2020distributed}. Distributed databases improve data access and processing because data is distributed to various network locations. Data access is faster, points of failure are less likely, and users have local control over their data \cite{abdelhafiz2020distributed}. \\

\noindent One widely adopted technique to improve the performance of distributed databases is sharding, which involves partitioning a large dataset from a single database node into smaller, more manageable segments called ``shards'' \cite{abdelhafiz2021sharding, bagui2015database}. These data shards or partitions are then distributed and stored across multiple geographically separated database servers. Sharding, exemplifying a shared-nothing architecture of distributed databases, focuses on scaling out by leveraging multiple database servers and utilizing their resources. Middleware-assisted sharding further enhances the scalability and performance of sharded distributed databases by managing the data distribution, load balancing, and fault tolerance through a dedicated middleware layer \cite{bernstein1996middleware}.  Moreover, it provides an additional layer of abstraction, simplifying the development and management of sharded distributed databases. \\

\noindent The motivation behind this thesis is to gain a deeper understanding of sharded distributed databases. Through the design and implementation of a middleware prototype, this thesis explores and evaluates the performance of a middleware-assisted sharded distributed database system using consistent hash-based sharding approach.

%***************************************************************************
\section{Research Questions and Objectives}
\label{RQ}
This thesis focuses on answering the following research questions to gain deeper insights into the topic:
\begin{itemize}
\item[] \textbf{RQ1:} How does the performance of sharded distributed databases compare to single centralized database in terms of throughput, response time, and data distribution, and what key factors influence these performance aspects in both approaches?

\item[] \textbf{RQ2:} How do different middleware configurations, such as the number of shards and database servers affect performance and scalability and how does consistent hashing with virtual nodes approach impact data distribution in sharded distributed databases?

\item[] \textbf{RQ3:} What are the challenges and potential bottlenecks in implementing middleware-assisted sharding in distributed databases, and how can they be addressed to ensure optimal scalability and performance?
\end{itemize}

\noindent To address these research questions, this thesis sets the following objectives:
\begin{itemize}
\item[] \textbf{Objective 1:} Design and implement a middleware-assisted sharded distributed database system with a focus on consistent hashing using virtual nodes approach.
\item[] \textbf{Objective 2:} Conduct a performance analysis and evaluation of the implemented middleware prototype by comparing it to a centralized database using performance metrics such as throughput, response time, and data distribution.
\item[] \textbf{Objective 3:}  Analyze the impact of different middleware configurations on performance and scalability and study the effects of consistent hashing on data distribution. Identify challenges and potential bottlenecks in implementing middleware-assisted sharding and propose strategies and future works to address them.
\end{itemize}


%***************************************************************************
\section{Thesis Outline}
This thesis focuses on the study and evaluation of middleware-assisted system of sharded distributed databases through the design and implementation of a prototype and is structured into seven chapters, each of which serves a specific purpose in achieving the research objectives. \\

\noindent \textbf{Chapter 1} is the introduction, which provides a brief overview for the motivation behind the research, the most important objectives, research questions followed by this section presenting the outline of the thesis providing a roadmap for the reader. \\

\noindent \textbf{Chapter 2} focuses on the background of distributed databases, including their types and architectures. The chapter specifically explores data distribution and sharding approaches and provides an overview of related works in the area, with a particular emphasis on sharding.\\

\noindent \textbf{Chapter 3} explores the design principles and considerations for implementing a middleware-assisted system of sharded distributed databases. This chapter presents the middleware’s design goals, functions and specifications, and its interacting modules and communication entities including underlying distributed databases and connected clients. The chapter also addresses the scope and limitations of the middleware’s design. \\

\noindent \textbf{Chapter 4} covers the implementation of the middleware-assisted system of sharded distributed databases. This chapter details the implementation process, including the infrastructures used to develop the middleware, the operations supported by the middleware, and concrete implementation of sharding in Redis databases. \\

\noindent \textbf{Chapter 5} evaluates the middleware’s performance using defined performance metrics of throughput and response time, custom workload, and custom benchmarking process. This chapter provides a comprehensive assessment and analysis of the middleware’s performance across various test scenarios, featuring different middleware configurations. Additionally, this chapter presents a brief analysis of data distribution across the associated Redis databases, emphasizing uniform data distribution. \\

\noindent \textbf{Chapter 6} provides a discussion of the research findings and their implications, including a critical evaluation of the research outcomes and a comparison with related works in the literature. \\

\noindent Finally, \textbf{Chapter 7} concludes this thesis by addressing all presented research questions, summarizing all evaluation results and outcomes, and highlighting implications for future work in the domain of sharded distributed databases.

